\chapter{The full research question}\label{ch:question}

The one-standoff restart was a starting point, not the goal. On 7 November you restated the actual question:

\begin{keyidea}[The research question]
Can a physics-based learned inverse estimate unknown standoff and reconstruct near-plane QDM magnetic fields
consistently across unseen geometries, fields of view and sampling scales?
\end{keyidea}

A paper on this would need to show:
\begin{enumerate}
\item joint standoff estimation and field reconstruction, with a distinct inference method;
\item generalisation to unseen geometries, fields of view, pixel spacings and noise levels;
\item how the method handles the ambiguity between source and height, for example with uncertainty or competing
hypotheses;
\item ideally, evidence on experimental QDM data.
\end{enumerate}

This chapter collects what we found when we looked at the question as a whole. Nothing in it was trained.

\section{An exact ambiguity}

Model the source as a thin current sheet at depth $h$ below the sensor. Every field component then has the form
\begin{equation}
B(k) = \tfrac{\mu_0}{2}\,e^{-kh}\,K(k)\,g(k),
\end{equation}
where $g$ describes the current and $K$ is a fixed factor that depends only on direction. Now move the sheet down by
$\delta$ and make it sharper by $e^{k\delta}$:
\begin{equation}
e^{-k(h+\delta)}\,K(k)\,\big(g(k)\,e^{k\delta}\big) = e^{-kh}\,K(k)\,g(k).
\end{equation}
\textbf{The measurement is identical.} A source at depth $h$ and a sharper source at depth $h+\delta$ cannot be told apart
from the field alone. Measuring $B_x$, $B_y$ and $B_z$ does not help, because all three share the same $e^{-kh}$.

\begin{keyidea}[What follows from it]
\begin{itemize}
\item \textbf{The height cannot come from the field alone.} Only a prior on what sources look like (sharp wire edges,
smooth currents, no ringing) can pick one depth.
\item \textbf{The two hypotheses predict different near-plane fields.} So a wrong height directly spoils the
reconstruction. Height and reconstruction are one problem.
\item \textbf{Where the prior is weak, the honest answer is a set of heights,} each with its own reconstructed field.
That is the direct answer to ``how do you handle the ambiguity''.
\end{itemize}
\end{keyidea}

\begin{careful}
``Standoff'' here means the \emph{effective} depth of the sheet. The thickness of the NV layer, the thickness of the
conductor and the optical blur all add to it. A paper should report the physical gap and the effective depth separately.
\end{careful}

\section{Can the spectrum tell the height?}

A tempting shortcut is to read the height off the slope of the field's power spectrum, since $|B|^2$ falls like
$e^{-2kh}$. We tried it on our COMSOL pairs.
\begin{itemize}
\item \textbf{With the fitting band chosen using the true height} (cheating): IC\_4um came out 0.2--0.6\um\ too high, but
L1\_C1 above 4\um\ came out 1--5\um\ too low; 10\um\ read as 5.2\um. The coil's own spectrum fights the exponential.
\item \textbf{With the band chosen from the data alone:} every height above 3\um\ read as 0.9--2\um, because COMSOL's
speckle is not white noise and leaks into the band.
\end{itemize}
Both failures are the ambiguity above in action: the answer depends on what the source looks like. The height has to
come from a source prior together with the reconstruction.

\section{What others have done}

\begin{center}
\small
\begin{tabularx}{\textwidth}{>{\raggedright\arraybackslash}p{4.2cm}X}
\toprule
Work & What it offers us \\
\midrule
NV practice (several papers) & Standoff fitted on a known feature: a wire profile, a film edge, a dipole. This is how
the real QDM could be calibrated. \\
Setescak et al.\ 2026 & Fits the sensor distance jointly with the magnetisation, using an energy prior. The closest prior
work, but for magnetisation, not learned, and with no uncertainty. \\
MAGIC-UNet, Reed et al.\ 2025 & A U-Net from field to current, trained over a range of standoffs that it is never told.
Chapters~\ref{ch:magic} and~\ref{ch:magic_results}. \\
Dubois et al.\ 2022 & An untrained network whose loss is the exact forward model. No training data, so no simulator gap. \\
USRNet, DPIR & One unrolled network for many blur kernels, scales and noise levels, given the kernel. \\
IKC, DAN & Alternate between estimating the blur and restoring the image. \\
BlindDPS, Fast Diffusion EM & A diffusion prior on the image, with the blur parameters estimated alongside. \\
\bottomrule
\end{tabularx}
\end{center}

\begin{keyidea}[The gap]
We found no learned method that estimates the standoff jointly with the reconstruction, reports competing heights, or is
tested across pixel spacing and field of view. The blind-deblurring literature has the machinery; nobody has applied it to
magnetic imaging. Our blur has only one unknown, the height, so a much more careful treatment is possible than in general
blind deblurring.
\end{keyidea}

\begin{careful}[Where these summaries come from]
Some entries were read in full; others come only from abstracts. Read the full papers before quoting any of them. The
survey with sources is \code{RESEARCH_BEYOND_ISTA.md}.
\end{careful}

\section{A candidate method, and the first test}

One design that fits the question:
\begin{enumerate}
\item For each candidate depth $h_j$ on a grid, run one shared reconstructor that is told $h_j$.
\item Score each pair $(h_j, x_j)$ by how well it fits the data and how plausible the source is.
\item Report the spread of supported heights, with one field for each. Give a single answer only when one height clearly
wins.
\item Refine the height around the best candidate.
\end{enumerate}
Every component of this is published (USRNet, DAN, BlindDPS); the combination and the treatment of the ambiguity would be
new. That keeps the professor's rule: published parts and baselines, with a clearly labelled contribution.

Before building any of it, one cheap check decides whether joint estimation can work at all. On COMSOL, fit a regularised
current at each candidate height and see whether the best fit lands on the true height. If no prior puts the minimum near
the truth, the height must be calibrated externally, and that is itself a result. This check has not been run yet.

\begin{defend}[If your professor asks: why not just estimate the height from the data?]
Because the data alone cannot fix it. A source at depth $h$ and a sharper one deeper down give exactly the same field. Our
spectral test on COMSOL confirmed it: the estimates were off by up to 5\um. Height can only come from a prior on sources,
and the method has to say how sure it is.
\end{defend}
